% !TeX root = ../../Book5.tex
% 纲要标题：### 17.1 Weyl 群
% 纲要位置：工作记录/计划纲要.md，第 4615 行。

\section{Weyl 群}
\label{b5:sec:1701}

从单个根反射走向它们生成的群，需要同时研究反射乘积和正根符号。交换引理将把最短表达、Weyl 室及支配代表联系起来。

% 纲要要点已在下文展开。

% 纲要细目（保留原定写作范围）：
%
% * 简单反射
% * Weyl 室与长度
% * 权格和根格
% * 用秩二交换计算和逐次反射说明长度与支配代表，用秩一自然表示区分根格与权格
%

\subsection{简单反射}
\label{b5:subsec:1701:01}

\begin{definition}\label{b5:def:weyl-group}
	设 \(E\) 为有限维实内积空间，\(\Phi\subseteq E\) 为有限约化结晶根系，\(s_\alpha\) 为根 \(\alpha\) 对应的正交反射。由这些反射生成的正交变换群称为 \(\Phi\) 的\term{Weyl 群}[Weyl group]，记为 \(W\)。固定简单根基 \(\Delta=\{\alpha_1,\ldots,\alpha_r\}\) 后，记 \(s_i=s_{\alpha_i}\)。
\end{definition}
\symidx{\(W,\Delta\)：根系的 Weyl 群与所选简单根基}

因为 \(W\) 置换有限集合 \(\Phi\)，而 \(\Phi\) 生成 \(E\)，故这个置换作用忠实，\(W\) 是有限群。\cref{b5:prop:simple-reflection-positive}还说明它由简单反射生成：每个根写成 \(w\alpha_i\)，其中 \(w\) 是简单反射的乘积，则 \(s_{w\alpha_i}=ws_iw^{-1}\)。

\begin{lemma}[删去一个反射]\label{b5:lem:weyl-exchange}
	设 \(\Phi\subseteq E\) 为有限约化结晶根系，选定正根集合及其简单根基 \(\{\alpha_1,\ldots,\alpha_r\}\)，记 \(s_i=s_{\alpha_i}\) 为相应简单反射。设 \(w=s_{i_1}\cdots s_{i_m}\)。若 \(w\alpha_j\) 为负根，则 \(ws_j\) 可由所列乘积删去某一个因子得到。
\end{lemma}
\begin{proof}
	从右向左依次把各反射作用于 \(\alpha_j\)。它起初为正，最终为负，故某一步首次把正根变成负根。 简单反射只把自己的简单根变为负根，因此存在 \(k\)，使
	\[
		v\alpha_j=\alpha_{i_k},\qquad v=s_{i_{k+1}}\cdots s_{i_m}.
	\]
	反射的共轭公式给出 \(vs_jv^{-1}=s_{i_k}\)，故
	\[
		ws_j =s_{i_1}\cdots s_{i_{k-1}}s_{i_k}vs_j =s_{i_1}\cdots s_{i_{k-1}}v.
	\]
	这正是删去第 \(k\) 个因子的乘积。
\end{proof}

在 \(A_2\) 中，简单反射满足 \(s_1\alpha_2=s_2\alpha_1=\alpha_1+\alpha_2\)。取 \(w=s_1s_2s_1\)，从右向左作用于 \(\alpha_2\)，得到
\[
	\alpha_2\xrightarrow{s_1}\alpha_1+\alpha_2 \xrightarrow{s_2}\alpha_1\xrightarrow{s_1}-\alpha_1.
\]
首次变号发生在最左侧的 \(s_1\)，所以交换引理删去原表达的第一个因子，给出
\[
	(s_1s_2s_1)s_2=s_2s_1.
\]
具体地，\(v=s_2s_1\) 满足 \(v\alpha_2=\alpha_1\)，从而 \(vs_2v^{-1}=s_1\)，这正是证明中的消去步骤。删去的因子不必位于原乘积的末端。再在右边乘 \(s_2\)，还得到 \(s_1s_2s_1=s_2s_1s_2\)。

\subsection{Weyl 室与长度}
\label{b5:subsec:1701:02}

全部根超平面 \(H_\alpha=\bigl\{\,x\in E\mid(x,\alpha)=0\,\bigr\}\) 把其补集分成若干连通分支，称为\term{Weyl 室}[Weyl chamber]。 选定简单根所对应的基本室为
\[
	C=\bigl\{\,x\in E\mid(x,\alpha_i)>0\;\text{对所有}\;i\,\bigr\}.
\]
其闭包 \(\overline C=\{x\in E\mid(x,\alpha_i)\ge0\;\text{对所有}\;i\}\) 称为闭基本室，也称闭支配室；属于它的向量称为支配向量。 基本室 \(C\) 确是一个连通分支：正根为简单根的非负整数和，所以 \(C\) 中所有正根的值都为正；它又是若干开半空间的交，因而凸且连通。任何不穿过根超平面的连续路径都不能改变这些符号。

\begin{theorem}\label{b5:thm:weyl-chambers-length}
	设 \(\Phi\subseteq E\) 为有限约化结晶根系，\(E\) 为实内积空间，\(W\) 为其 Weyl 群。选定正根集合 \(\Phi^+\) 及简单根基 \(\{\alpha_1,\ldots,\alpha_r\}\)，令 \(C=\{v\in E\mid(v,\alpha_i)>0\;\text{对所有}\;i\}\) 为基本 Weyl 室，\(\overline C\) 为其闭包。对 \(w\in W\)，令 \(\ell(w)\) 为其最短简单反射乘积的因子数。则
	\begin{equation}\label{b5:eq:weyl-length-inversions}
		\ell(w)=\#\{\,\alpha\in\Phi^+\mid w\alpha\in-\Phi^+\,\}.
	\end{equation}
	此外，\(W\) 在 Weyl 室上单纯传递，每个 \(W\)-轨道与闭基本室 \(\overline C\) 相交于唯一一点。
\end{theorem}
\begin{proof}
	将右边记为 \(N(w)\)。因 \(s_i\) 置换除 \(\alpha_i\) 外的正根，有
	\[
		N(ws_i)=
		\begin{cases}
			N(w)+1,&w\alpha_i>0,\\
			N(w)-1,&w\alpha_i<0.
		\end{cases}
	\]
	对一个最短表达 \(w=us_i\)，其前缀也是最短表达。若 \(u\alpha_i<0\)，\cref{b5:lem:weyl-exchange}会把 \(us_i\) 写成比当前表达少两个因子的乘积，矛盾。因此 \(u\alpha_i>0\)，按长度归纳得到 \(N(w)=\ell(w)\)。

	取 \(c\in C\)。对任意 \(x\in E\)，在有限轨道 \(Wx\) 中选取与 \(c\) 内积最大的 \(y\)。若 \((y,\alpha_i)<0\)，则
	\[
		(s_i y,c)-(y,c) =-\frac{2(y,\alpha_i)(\alpha_i,c)}{(\alpha_i,\alpha_i)}>0,
	\]
	矛盾。所以 \(y\in\overline C\)。若 \(x\) 不在任何根超平面上，则 \(y=wx\) 也不在其中，因为 \(W\) 置换根超平面。因而所有 \((y,\alpha_i)\ge0\) 都不能为零，实际均为正，即 \(y\in C\)。这说明各室均为某个 \(wC\)。 若 \(wC=C\)，则 \(w\) 保持正根，故 \(N(w)=0\)，从而 \(w=1\)。

	还要处理闭室的边界。若 \(w=s_{i_1}\cdots s_{i_m}\) 最短，则每个 \(\beta_k=s_{i_1}\cdots s_{i_{k-1}}\alpha_{i_k}\) 为正根；否则交换引理使相应前缀变短。相邻前缀之差相加，得到算子恒等式
	\[
		1-s_{i_1}\cdots s_{i_m} =\sum_{k=1}^m s_{i_1}\cdots s_{i_{k-1}}(1-s_{i_k}),
	\]
	其中空乘积为恒等算子。将它作用于任意支配向量 \(\mu\in\overline C\)，再用反射公式，得到
	\begin{equation}\label{b5:eq:dominant-weyl-difference}
		\mu-w\mu =\sum_{k=1}^m \frac{2(\mu,\alpha_{i_k})}{(\alpha_{i_k},\alpha_{i_k})}\,\beta_k.
	\end{equation}
	所以 \(\mu-w\mu\) 是简单根的非负实线性组合。若 \(w\mu\) 也支配，反过来用 \(w^{-1}\)，则 \(w\mu-\mu\) 也是这样的组合。简单根线性无关，迫使二者都为零。因此闭室中的代表唯一。
\end{proof}

在 \(A_{n-1}\) 的标准坐标中，\(W=S_n\) 置换坐标；一个基本室由 \(x_1>x_2>\cdots>x_n\)、\(\sum x_i=0\) 给出。 长度公式便成为排列的逆序数公式。坐标相等时，向量落在室的边界上，其稳定子群可以非平凡，这与室内部的单纯传递性并不冲突。

\ProseParagraphBreak

在 \(A_2\) 中，将简单根画成 \(\alpha_1=(1,0)\)、\(\alpha_2=(-1/2,\sqrt3/2)\)。两根之间的角是 \(120^\circ\)，基本室却由它们的垂直墙围成，角为 \(60^\circ\)。从 \(C\) 依次走到 \(s_1C,s_1s_2C,s_1s_2s_1C\)，穿过的墙分别垂直于 \(\alpha_1,\alpha_1+\alpha_2,\alpha_2\)。终点是相反室，三个正根全都变号，长度公式因而给出 \(\ell(s_1s_2s_1)=3\)。这些墙与根的对应画在\cref{b5:fig:a2-chambers-lattices}中；墙上的基本权将在下面给出坐标。

\subsection{权格和根格}
\label{b5:subsec:1701:03}

\begin{definition}\label{b5:def:weight-lattice}
	设 \(E\) 为有限维实内积空间，\(\Phi\subseteq E\) 为有限约化结晶根系，固定简单根基 \(\{\alpha_1,\ldots,\alpha_r\}\)，记 \(\overline C=\{v\in E\mid(v,\alpha_i)\ge0\;\text{对所有}\;i\}\) 为闭基本室。 令 \(\alpha^\vee=2\alpha/(\alpha,\alpha)\)，并记 \(\langle\lambda,\alpha^\vee\rangle=2(\lambda,\alpha)/(\alpha,\alpha)\)。 两个格
	\[
		Q=\sum_i\mathbb Z\alpha_i,\qquad P=\bigl\{\,\lambda\in E\mid \langle\lambda,\alpha^\vee\rangle\in\mathbb Z \;\text{对所有}\;\alpha\in\Phi\,\bigr\}
	\]
	分别称为\term{根格}[root lattice]和\term{权格}[weight lattice]。 满足 \(\langle\omega_i,\alpha_j^\vee\rangle=\delta_{ij}\) 的向量 \(\omega_i\) 称为\term{基本权}[fundamental weight]。 属于 \(P\cap\overline C\) 的向量称为\term{支配整权}[dominant integral weight]，其集合记为 \(P^+\)。
\end{definition}
\symidx{\(\alpha^\vee,\langle\lambda,\alpha^\vee\rangle\)：余根及权与余根的配对}
\symidx{\(Q,P\)：根格与权格}
\symidx{\(\omega_i,P^+\)：基本权与支配整权集合}

\begin{proposition}\label{b5:prop:weight-fundamental-basis}
	设 \(\Phi\subseteq E\) 为有限约化结晶根系，\(\{\alpha_1,\ldots,\alpha_r\}\) 为简单根基，\(\omega_i\) 为相应基本权，即 \(\langle\omega_i,\alpha_j^\vee\rangle=\delta_{ij}\)。记 \(P\) 为权格，\(Q\) 为根格，\(P^+\) 为支配整权集合。则
	\[
		P=\bigoplus_i\mathbb Z\omega_i,\qquad P^+=\sum_i\mathbb Z_{\geq0}\omega_i,\qquad Q\subseteq P,
	\]
	且 \(P/Q\) 有限。检验一个向量是否属于 \(P\)，只须检验它与简单余根的配对是否全为整数。
\end{proposition}
\begin{proof}
	余根集合本身也是根系：反射与原根系相同，结晶整数交换了原来的两个根参数。 若正根 \(\beta=\sum_i c_i\alpha_i\)、\(c_i\ge0\)，则
	\[
		\beta^\vee=\frac{2\beta}{(\beta,\beta)} =\sum_i c_i\frac{(\alpha_i,\alpha_i)}{(\beta,\beta)}\alpha_i^\vee.
	\]
	系数均非负，所以每个正根的余根都是简单余根的非负实组合。若一个简单余根是两个正余根之和，比较在简单余根基下的非负实系数，两个加数都只能在该简单余根的射线上；约化性使它们都等于该简单余根，矛盾。 因此全部简单余根都是余根系的简单根；它们已有 \(r\) 个，由\cref{b5:thm:simple-roots}，正好组成余根系的简单根基。 故每个余根都是简单余根的整数和，检验简化随之成立。 基本权是简单余根的对偶基，所以所列两种展开立即成立。结晶条件给出 \(Q\subseteq P\)。 令 \(A=(a_{ij})\)、\(a_{ij}=\langle\alpha_j,\alpha_i^\vee\rangle\)。由基本权的对偶关系， \(\alpha_j=\sum_i a_{ij}\omega_i\)，故在基本权坐标下 \(P=\mathbb Z^r\)、\(Q=A\mathbb Z^r\)。 矩阵 \(A\) 可逆且为整数矩阵，\(\operatorname{adj}(A)\) 也是整数矩阵；恒等式 \(A\operatorname{adj}(A)=(\det A)I\) 给出 \((\det A)\omega_j\in Q\) 对每个 \(j\) 成立。因此有限群 \(P/(|\det A|P)\) 满射到 \(P/Q\)，后者也有限。
\end{proof}

根格与权格的差别在秩一已经出现。对 \(A_1\) 根系 \(\Phi=\{\alpha,-\alpha\}\)，基本权为 \(\omega=\dfrac{\alpha}{2}\)，因为 \(\langle\omega,\alpha^\vee\rangle=1\)。于是
\[
	\begin{aligned}
		Q&=\mathbb Z\alpha=2\mathbb Z\omega,\qquad P=\mathbb Z\omega,\\
		P/Q&\cong\mathbb Z/2\mathbb Z,\qquad P^+=\mathbb Z_{\ge0}\omega.
	\end{aligned}
\]

\ProseParagraphBreak

这个额外的半根在表示中确实出现。对 \(\mathfrak{sl}_2(\mathbb C)\)，取标准基 \(e,f,h\)，正根满足 \(\alpha(h)=2\)；\cref{b5:thm:root-sl2}中的余根元 \(h_\alpha\) 在这里就是 \(h\)。自然二维模 \(L_1\) 的最高向量 \(v\) 满足 \(ev=0\)、\(hv=v\)，故其最高权 \(\lambda\) 满足 \(\lambda(h)=1\)，即 \(\lambda=\omega=\dfrac{\alpha}{2}\notin Q\)。更一般地，\cref{b5:prop:sl2-model,b5:thm:sl2-classification}中的 \(L_m\) 以 \(hv_0=mv_0\) 给出最高权 \(m\omega\)，所以这些有限维不可约模的最高权恰好取遍 \(P^+\)。根格记录根的整数和，权格则容纳了这些并非都在根格中的最高权。

\ProseParagraphBreak

正简单根实锥与闭支配室使用两套非负条件。若 \(K=\sum_i\mathbb R_{\ge0}\alpha_i\)，则
\[
	\overline C=\{\,x\in E\mid(x,y)\ge0\;\text{对所有}\;y\in K\,\},
\]
所以闭支配室是 \(K\) 的内积对偶锥。在上述 \(A_2\) 坐标中，配对方程给出
\[
	\omega_1=(1/2,\sqrt3/6),\qquad \omega_2=(0,\sqrt3/3), \qquad \alpha_1=2\omega_1-\omega_2,\quad \alpha_2=-\omega_1+2\omega_2.
\]
因此正根 \(\alpha_2\) 不是支配权，而 \(\omega_1\) 是支配整权却不在根格中：由两式解得 \(\omega_1=(2\alpha_1+\alpha_2)/3\)，其简单根系数不是整数。又有 \((\omega_1,\alpha_2)=0\)，故 \(s_2\omega_1=\omega_1\)。图中的格点及墙同时展示了这些区别；边界点有稳定子群，不影响开室集合上的单纯传递作用。

% !TeX root = ../Book5.tex
\begin{figure}[htbp]
	\centering
	\begin{tikzpicture}[x=2.7cm,y=2.7cm,every node/.style={font=\small}]
		\fill[yellow!12] (0,0)--(0:1.18) arc[start angle=0,end angle=120,radius=1.18]--cycle;
		\fill[AlgebraBlue!12] (0,0)--(30:1.18) arc[start angle=30,end angle=90,radius=1.18]--cycle;
		\foreach \m in {-3,...,3} {
			\foreach \n in {-3,...,3} {
				\pgfmathsetmacro{\px}{\m/2}
				\pgfmathsetmacro{\py}{sqrt(3)*(\m/6+\n/3)}
				\pgfmathtruncatemacro{\inside}{(\px*\px+\py*\py<2.35)?1:0}
				\ifnum\inside=1
					\fill[gray!45] (\px,\py) circle[radius=.6pt];
					\pgfmathtruncatemacro{\rootpoint}{mod(\m+2*\n,3)}
					\ifnum\rootpoint=0
						\fill[black] (\px,\py) circle[radius=1.1pt];
					\fi
				\fi
			}
		}
		\foreach \angle in {30,90,150} {
			\draw[dashed,gray] (\angle:1.60)--(\angle+180:1.60);
		}
		\foreach \angle/\name in {0/\alpha_1,60/{\alpha_1+\alpha_2},120/\alpha_2,
			180/{-\alpha_1},240/{-\alpha_1-\alpha_2},300/{-\alpha_2}} {
			\draw[->,line width=.7pt] (0,0)--(\angle:1);
			\node[fill=white,inner sep=1pt] at (\angle:1.17) {\(\name\)};
		}
		\draw[->,AlgebraBlue,line width=.9pt] (0,0)--(.5,.288675);
		\draw[->,AlgebraBlue,line width=.9pt] (0,0)--(0,.577350);
		\node[anchor=north west,text=AlgebraBlue] at (.5,.288675) {\(\omega_1\)};
		\node[anchor=east,text=AlgebraBlue] at (0,.577350) {\(\omega_2\)};
		\node[anchor=north east] at (0,0) {\(0\)};
		\foreach \angle in {60,120,180} {
			\draw[->,AlgebraBlue,line width=.7pt]
				(\angle:1.48) arc[start angle=\angle,end angle=\angle+55,radius=1.48];
		}
		\node at (60:1.73) {\(C\)};
		\node at (120:1.73) {\(s_1C\)};
		\node at (180:1.85) {\(s_1s_2C\)};
		\node at (240:1.75) {\(s_1s_2s_1C=-C\)};
	\end{tikzpicture}
	\caption[\(A_2\) 的室与两种格]{\(A_2\) 的六根、三条墙及室链。 淡黄色扇形是正简单根实锥，蓝色扇形是基本室；二者的边界不同。 灰点表示权格，黑点表示根格。\(\omega_1\) 位于垂直于 \(\alpha_2\) 的墙上，故被 \(s_2\) 固定。 外侧箭头给出到相反室的三步最短室链。}
	\label{b5:fig:a2-chambers-lattices}
\end{figure}


\begin{proposition}\label{b5:prop:weyl-dominant-order}
	设 \(\Phi\subseteq E\) 为有限约化结晶根系，\(\{\alpha_1,\ldots,\alpha_r\}\) 为所选简单根基，\(W\) 为 Weyl 群，\(P\) 为权格，\(P^+\) 为支配整权集合。每个 \(\mu\in P\) 的 \(W\) 轨道都有唯一支配整权。若 \(\mu\in P^+\)、\(w\in W\)，则 \(\mu-w\mu\in\sum_i\mathbb Z_{\geq0}\alpha_i\)。
\end{proposition}
\begin{proof}
	反射公式说明 \(W\) 保持 \(P\)，故闭室代表仍为整权。存在唯一性已经由\cref{b5:thm:weyl-chambers-length}证明。公式\eqref{b5:eq:dominant-weyl-difference}中的系数此时都是非负整数，而 \(\beta_k\) 为正根，故给出后一结论。
\end{proof}

\begin{proposition}\label{b5:prop:rho-coroot}
	设 \(\Phi\subseteq E\) 为有限约化结晶根系，\(\Phi^+\) 为所选正根集合，\(\{\alpha_1,\ldots,\alpha_r\}\) 为相应简单根基，\(\omega_i\) 为基本权，\(\alpha_i^\vee\) 为简单余根。令 \(\rho=\dfrac12\sum_{\alpha\in\Phi^+}\alpha\)。则 \(\rho=\sum_i\omega_i\)，特别地 \(\langle\rho,\alpha_i^\vee\rangle=1\)。
\end{proposition}
\symidx{\(\rho=\frac12\sum_{\alpha\in\Phi^+}\alpha\)：正根和的一半}
\begin{proof}
	\(s_i\) 置换除 \(\alpha_i\) 外的正根，故 \(s_i\rho=\rho-\alpha_i\)。 与反射公式比较便得每个配对等于 \(1\)，再用基本权基展开。
\end{proof}

支配代表也可以逐步求出。对任意 \(\mu\in P\)，从 \(\mu_0=\mu\) 开始；若当前有 \(\langle\mu_k,\alpha_i^\vee\rangle<0\)，便任选这样的 \(i\)，令
\[
	\mu_{k+1}=s_i\mu_k =\mu_k-\langle\mu_k,\alpha_i^\vee\rangle\alpha_i.
\]
固定 \(c\in C\)，每一步都满足
\[
	(\mu_{k+1},c)-(\mu_k,c) =-\langle\mu_k,\alpha_i^\vee\rangle(\alpha_i,c)>0.
\]
所有 \(\mu_k\) 都在有限轨道 \(W\mu\) 中，严格增加使任何点都不能重复，故至多经过 \(\#(W\mu)-1\) 次反射便终止。终止时所有简单余根配对非负，且反射始终保持权格，所以所得点属于 \(P^+\)；\cref{b5:prop:weyl-dominant-order}保证它是唯一的支配代表。

\ProseParagraphBreak

例如在 \(A_2\) 中，\(\rho=\alpha_1+\alpha_2=\omega_1+\omega_2\)。从 \(-\rho\) 出发，先选 \(s_1\)，再依次选 \(s_2,s_1\)，有
\[
	-\rho\xrightarrow{s_1}-\alpha_2 \xrightarrow{s_2}\alpha_2\xrightarrow{s_1}\rho.
\]
各点与 \((\alpha_1^\vee,\alpha_2^\vee)\) 的配对依次为
\[
	(-1,-1),\qquad(1,-2),\qquad(-1,2),\qquad(1,1).
\]
每次选用的反射都对应一个负配对；最后两个配对均为正，过程到达基本室内部。
